\section{Riemann and Darboux Integration}

\subsection{Foundational Definitions}

\begin{enumerate}
  \item What is a partition of an interval $[a,b]$?
  \item Define the norm (mesh size) of a partition.
  \item What is a tagged partition?
  \item Define a Riemann sum.
  \item State the definition of the Riemann integral using Riemann sums.
  \item Define the upper Darboux sum.
  \item Define the lower Darboux sum.
  \item What are the quantities
    \[
      M_i = \sup_{x \in [x_{i-1},x_i]} f(x)
    \]
    and
    \[
      m_i = \inf_{x \in [x_{i-1},x_i]} f(x)?
    \]
  \item Define the upper Darboux integral.
  \item Define the lower Darboux integral.
  \item State the definition of Darboux integrability.
  \item Why must a function be bounded to be Riemann integrable?
  \item Why must a function be bounded to be Darboux integrable?
\end{enumerate}

\subsection{Conceptual Understanding}

\begin{enumerate}[resume]
  \item What is the geometric intuition behind a Riemann sum?
  \item What is the geometric intuition behind upper and lower Darboux sums?
  \item Why does refining a partition generally improve approximation quality?
  \item What happens to lower Darboux sums when a partition is refined?
  \item What happens to upper Darboux sums when a partition is refined?
  \item Why is
    \[
      L(f,P) \le U(f,P)
    \]
    always true?
  \item Explain why
    \[
      \sup_P L(f,P) \le \inf_P U(f,P).
    \]
  \item What does it mean geometrically when the upper and lower integrals are equal?
  \item Why are Darboux sums often considered technically simpler than Riemann sums?
  \item Why are Riemann sums closer to the historical intuition of integration?
\end{enumerate}

\subsection{Comparisons Between Riemann and Darboux Integration}

\begin{enumerate}[resume]
  \item What is the main difference between a Riemann sum and a Darboux sum?
  \item In a Riemann sum, how are sample points chosen?
  \item In a Darboux sum, why are sample points unnecessary?
  \item Which theory uses arbitrary sample points explicitly?
  \item Which theory uses suprema and infima on subintervals?
  \item What role does the partition norm play in Riemann integration?
  \item What role does partition refinement play in Darboux integration?
  \item State the theorem relating Riemann integrability and Darboux integrability.
  \item Why are the two theories equivalent for bounded functions on closed intervals?
  \item Can two different tagged partitions with the same underlying partition produce different Riemann sums?
  \item Can two different Darboux sums exist for the same partition?
  \item How do upper and lower Darboux sums bound all possible Riemann sums?
  \item Explain how Darboux integration avoids explicit limits of sums.
  \item Explain how Riemann integration encodes convergence directly through limits.
\end{enumerate}

\subsection{Theorems and Properties}

\begin{enumerate}[resume]
  \item State the linearity property of the integral.
  \item State the additivity over intervals property.
  \item State the monotonicity property of the integral.
  \item If $f$ is continuous on $[a,b]$, why is it integrable?
  \item Why is every monotone function on $[a,b]$ integrable?
  \item State a criterion for integrability involving discontinuities.
  \item What is the relationship between measure-zero discontinuity sets and integrability?
  \item Why is the Dirichlet function not Riemann integrable?
  \item Why is the characteristic function of the rationals discontinuous everywhere?
  \item Give an example of a function that is discontinuous but still Riemann integrable.
  \item Why is a function with finitely many discontinuities integrable?
  \item What happens to the integral if two functions differ only on finitely many points?
  \item What happens if two functions differ only on a measure-zero set?
\end{enumerate}

\subsection{Technical Manipulation}

\begin{enumerate}[resume]
  \item Compute the lower and upper Darboux sums for a constant function.
  \item Compute the lower and upper Darboux sums for $f(x)=x$ on $[0,1]$ using a uniform partition.
  \item Derive the Riemann sum for $f(x)=x^2$ on $[0,1]$.
  \item Show directly that
    \[
      \int_0^1 x\,dx = \frac12
    \]
    using Riemann sums.
  \item Show directly that
    \[
      \int_0^1 x\,dx = \frac12
    \]
    using Darboux sums.
  \item For a uniform partition, how do you compute
    \[
      \Delta x?
    \]
  \item Why do continuous functions attain maxima and minima on closed intervals?
  \item Why is compactness important in Darboux integration?
  \item What is the difference between taking a supremum over all lower sums and taking a limit?
  \item Why does the existence of arbitrarily close upper and lower sums imply integrability?
\end{enumerate}

\subsection{Interpreting the Integral}

\begin{enumerate}[resume]
  \item What does the expression
    \[
      \sum_i f(x_i)\Delta x_i
    \]
    represent before assigning any geometric or physical interpretation to it?
  \item Why is the integral fundamentally an accumulation operator rather than specifically an area operator?
  \item In the area interpretation of the integral, what do the quantities
    \[
      f(x_i), \quad \Delta x_i, \quad f(x_i)\Delta x_i
    \]
    represent geometrically?
  \item Why does
    \[
      f(x_i)\Delta x_i
    \]
    approximate the area of a thin rectangle?
  \item What happens to the Riemann sum as
    \[
      \max \Delta x_i \to 0?
    \]
  \item Why does the definite integral emerge as a limit of Riemann sums?
  \item If
    \[
      g'(x)=f(x),
    \]
    why does
    \[
      f(x)\Delta x
    \]
    approximate the small change
    \[
      \Delta g?
    \]
  \item What is the conceptual meaning of
    \[
      dg = g'(x)\,dx?
    \]
  \item Why can the same expression
    \[
      f(x)\Delta x
    \]
    be interpreted both as rectangle area and as accumulated change?
  \item What determines the interpretation of the integrand \(f(x)\)?
  \item Why is the equality
    \[
      \int_a^b g'(x)\,dx = g(b)-g(a)
    \]
    an accumulation statement?
  \item In the Fundamental Theorem of Calculus, what exactly is being accumulated?
  \item Why does integrating a derivative recover net change?
  \item How does the telescoping intuition appear in
    \[
      \sum \Delta g_i?
    \]
  \item Define
    \[
      A(x)=\int_a^x f(t)\,dt.
    \]
    What does \(A(x)\) represent geometrically?
  \item Why does increasing \(x\) by a small amount \(h\) add approximately
    \[
      f(x)h
    \]
    to the accumulated area?
  \item Use the difference quotient to explain why
    \[
      A'(x)=f(x).
    \]
  \item Why does the identity
    \[
      A'(x)=f(x)
    \]
    connect the area interpretation to the antiderivative interpretation?
  \item Why is the Fundamental Theorem of Calculus fundamentally a statement that differentiation and integration are inverse operations?
  \item How does the notion of ``local rate'' connect to ``global accumulation''?
  \item What is the difference between:
    \begin{itemize}
      \item a function value,
      \item a derivative,
      \item a differential,
      \item an accumulated quantity?
    \end{itemize}
  \item Why is
    \[
      \int_a^b f(x)\,dx
    \]
    not intrinsically an area?
  \item Give examples where the integral represents:
    \begin{itemize}
      \item distance,
      \item mass,
      \item charge,
      \item probability,
      \item energy.
    \end{itemize}
  \item In each example above, what plays the role of:
    \begin{itemize}
      \item density/rate,
      \item small measure,
      \item accumulated total?
    \end{itemize}
  \item Explain the pattern:
    \[
      \text{density} \times \text{small measure}
      \to
      \text{accumulated quantity}.
    \]
  \item Why is the rectangle picture only one geometric realization of integration?
  \item What is the conceptual difference between:
    \[
      f(x)
      \quad \text{and} \quad
      f(x)\,dx?
    \]
  \item Why is
    \[
      dx
    \]
    best understood as ``a tiny increment of the independent variable''?
  \item Why is the statement
    \[
      dg=g'(x)\,dx
    \]
    more fundamental than the rectangle interpretation?
  \item What does the Fundamental Theorem of Calculus reveal about the relationship between local structure and global structure?
\end{enumerate}
